\section{The reporting funnel: where antisemitic incidents go missing}\label{sec:funnel}

The police count is the end of a sequence of losses: an incident must reach a recording channel, then the police; the police must record it as motivated by antisemitic bias; and the agency must submit the record to the FBI. Each stage loses incidents for different reasons and calls for a different remedy. This section formalises the sequence as a multiplicative identity (\cref{prop:funnel}), states which data identify which stage, and attributes differences in recorded rates between jurisdictions to the stages by a Shapley decomposition (\cref{thm:shapley}); \cref{fig:funnel} depicts it. In this section $U_d$ denotes \emph{all} in-scope incidents in domain $d$, whether or not they reach a channel, and $N$ their number; the population of \cref{sec:framework,sec:identification,sec:linkage,sec:models} is the subset $U_d^V$ defined below, of size $N_V$.

\subsection{The funnel identity}\label{sec:funnel-identity}

\begin{definition}[Funnel stages]\label{def:funnel}
A \emph{recording channel} is the police or an organisation that maintains one of the lists. For an incident $i\in U_d$ let $V_i=1$ if the incident reaches at least one recording channel, either through a report by the victim, a witness or another third party ($V_i^{\mathrm{rep}}=1$) or through the channel's own detection ($V_i^{\mathrm{det}}=1$), so that $V_i=\max(V_i^{\mathrm{rep}},V_i^{\mathrm{det}})$. Let $P_i=1$ if the incident reaches the police, by report or detection; $C_i=1$ if the police record it as motivated by antisemitic bias; and $S_i=1$ if the agency submits that record to the FBI's hate-crime collection. The stages are nested, $S_i\le C_i\le P_i\le V_i$, and $\nFBI=\sum_{i\in U_d}S_i$ is the FBI count. When list~1 is the FBI submission, $\nFBI=n_{1+}$; when list~1 is a local police list observed before submission, as in the pilot of \cref{sec:plan}, $n_{1+}=\sum_iC_i$ and the submission stage drops out. Let $I$ be uniform on $U_d$, independent of the indicators, and define the stage rates
\[
\begin{aligned}
  \rvic&=\Pr(V_I=1), &\qquad \rpol&=\Pr(P_I=1\mid V_I=1),\\
  \rcls&=\Pr(C_I=1\mid P_I=1), &\qquad \rsub&=\Pr(S_I=1\mid C_I=1).
\end{aligned}
\]
Probabilities are taken over the random incident and the reporting process, given $U_d$.
\end{definition}

The first stage is reaching a recording channel, not disclosure to anyone: an incident that the victim tells only to family has $V_i=0$ and is on no list. Where channels learn of incidents only through reports, $\rvic=\Pr(V^{\circ}_I=1)\Pr(V_I=1\mid V^{\circ}_I=1)$, with $V^{\circ}_i=1$ if the incident is disclosed to anyone; a survey that asks both questions identifies both factors, which respond to victim support and to the accessibility of channels respectively. The rates are \emph{aggregate}: $\rpol$ is the share that reaches the police among incidents that reach some channel, whatever their type. External sources observe a stage only on a subpopulation selected by their own mechanism. Let $W_i=1$ if incident $i$ is within the scope of a victim survey and reported by a sampled respondent; $G_i=1$ if $i$ is recorded by an NGO or community list ($k\ge2$) whose record states that the incident was reported to the police; and $A_i=1$ if the agency with jurisdiction over $i$ submits hate-crime data for the period, with $c_d$ the share of the domain's population served by such agencies.

\begin{assumption}[Source ignorability]\label{asm:funnel}
\leavevmode
\begin{enumerate}[({F}1)]
  \item\label{asm:funnel-survey} \emph{Survey.} $W_I\perp V_I$, and $W_I\perp P_I\mid V_I=1$. Within the survey frame, channels learn of incidents only through reports: $V^{\mathrm{det}}_I=0$ whenever $W_I=1$.
  \item\label{asm:funnel-overlap} \emph{Overlap.} $G_I=1$ implies $P_I=1$; $G_I\perp C_I\mid P_I=1$; and $G_I\perp S_I\mid C_I=1$.
  \item\label{asm:funnel-part} \emph{Participation.} $S_i=C_iA_i$ for every $i$; $A_I\perp C_I$; and $\Pr(A_I=1)=c_d$.
  \item\label{asm:funnel-mse} \emph{MSE target.} The capture model of \cref{sec:identification} or \cref{sec:models}, including its dependence restriction, holds on the incidents that reach a channel, $U_d^{V}=\{i:V_i=1\}$, with size $N_V=|U_d^{V}|$. (A list can record only incidents known to its organisation, so $\Pr(\h_i\neq\zero\mid V_i=0)=0$ by construction.)
\end{enumerate}
\end{assumption}

\begin{proposition}[Funnel identity]\label{prop:funnel}
\leavevmode
\begin{enumerate}[(a)]
  \item \emph{Accounting.} $\E[\nFBI\mid U_d]=N\,\rvic\,\rpol\,\rcls\,\rsub$, with no assumption beyond nesting. The realised analogue $\nFBI=N\cdot\frac{n_V}{N}\cdot\frac{n_P}{n_V}\cdot\frac{n_C}{n_P}\cdot\frac{\nFBI}{n_C}$, with $n_V=\sum_iV_i$, $n_P=\sum_iP_i$ and $n_C=\sum_iC_i$, holds whenever these counts are positive.
  \item \emph{Heterogeneity.} Suppose the stage probabilities of incident $i$ depend on $i$ only through a type $X_i$, such as a latent class (\cref{sec:models-lcmcr}) or a covariate profile (\cref{sec:models-cov}). Write $\rvic(x)=\Pr(V_i=1\mid X_i=x)$, $\rpol(x)=\Pr(P_i=1\mid V_i=1,X_i=x)$, and similarly $\rcls(x)$ and $\rsub(x)$. Let $X\sim G_d$, the empirical type distribution of $U_d$. Then each aggregate rate is a selection-weighted mean, for example $\rpol=\E[\rvic(X)\rpol(X)]/\E[\rvic(X)]$, and $\E[\nFBI\mid U_d]=N\,\E[\rvic(X)\rpol(X)\rcls(X)\rsub(X)]$. The factorisation into \emph{unweighted} type averages,
  \begin{equation}\label{eq:funnel-unweighted}
    \E[\nFBI\mid U_d]=N\,\E[\rvic(X)]\,\E[\rpol(X)]\,\E[\rcls(X)]\,\E[\rsub(X)],
  \end{equation}
  holds if the four random variables $\rho_{\mathrm{stage}}(X)$ are mutually independent under $G_d$, in particular if at most one stage varies with type. If $X$ is real-valued and all four stage probabilities are nondecreasing in $X$ (a common severity gradient), then the left side of \eqref{eq:funnel-unweighted} is at least the right side.
  \item \emph{Identification by source.} Under \ref{asm:funnel-survey}, $\rvic=\Pr(V_I=1\mid W_I=1)$ and $\rpol=\Pr(P_I=1\mid V_I=1,W_I=1)$. Under \ref{asm:funnel-overlap}, $\Pr(S_I=1\mid G_I=1)=\rcls\,\rsub$. Under \ref{asm:funnel-part}, $\rsub=c_d$. Under \ref{asm:funnel-mse}, MSE estimates $N_V$, and $\E[N_V\mid U_d]=N\rvic$, so that $\E[\nFBI\mid U_d]=\E[N_V\mid U_d]\,\rpol\,\rcls\,\rsub$.
\end{enumerate}
\end{proposition}

\begin{proof}
(a) By nesting, $\{S_I=1\}=\{V_I=1,P_I=1,C_I=1,S_I=1\}$, so the chain rule gives $\Pr(S_I=1)=\rvic\rpol\rcls\rsub$. Moreover, $\E[\nFBI\mid U_d]=\sum_i\Pr(S_i=1)=N\Pr(S_I=1)$. The realised identity telescopes.
(b) Conditioning on the type, $\Pr(V_I=1,P_I=1)=\E[\rvic(X)\rpol(X)]$ and $\Pr(V_I=1)=\E[\rvic(X)]$. Their ratio is $\rpol$, and the other stages are handled in the same way. The same argument gives $\Pr(S_I=1)=\E[\prod_{\mathrm{stage}}\rho_{\mathrm{stage}}(X)]$. Under mutual independence the expectation of the product is the product of the expectations. For the inequality, let $f,g\ge0$ be nondecreasing. Then $\E[f(X)g(X)]\ge\E f(X)\,\E g(X)$ by Chebyshev's association inequality, which follows from $\E[\{f(X)-f(X')\}\{g(X)-g(X')\}]\ge0$ for an independent copy $X'$. A product of nonnegative nondecreasing functions is nonnegative and nondecreasing, so induction over the four stages gives $\E[\prod\rho_{\mathrm{stage}}(X)]\ge\prod\E[\rho_{\mathrm{stage}}(X)]$.
(c) Under \ref{asm:funnel-survey}, $\Pr(V_I=1\mid W_I=1)=\Pr(V_I=1)$ and $\Pr(P_I=1\mid V_I=1,W_I=1)=\Pr(P_I=1\mid V_I=1)$, and within the survey frame $V_I=V_I^{\mathrm{rep}}$, which the survey observes. Under \ref{asm:funnel-overlap}, nesting and $G_I\Rightarrow P_I$ give $\Pr(S_I=1\mid G_I=1)=\Pr(C_I=1\mid P_I=1,G_I=1)\Pr(S_I=1\mid C_I=1,G_I=1)=\rcls\rsub$. Under \ref{asm:funnel-part}, $\rsub=\Pr(A_I=1\mid C_I=1)=\Pr(A_I=1)=c_d$. Under \ref{asm:funnel-mse}, every recorded incident lies in $U_d^V$ by construction, and the capture model holds on $U_d^V$, so the MSE estimand is $N_V$. Moreover $\E[N_V\mid U_d]=\sum_i\Pr(V_i=1)=N\rvic$, and the last display follows from (a).
\end{proof}

Part (a) is an accounting identity; the content lies in (b) and (c). Part (b) warns against assembling the funnel from rates measured on different type mixtures: under a common severity gradient, the product of unweighted stage averages understates the share of incidents that reach the FBI, so dividing $\nFBI$ by it overstates $N$. Part (c) says what each source delivers (\cref{tab:funnel-sources}). With all four sources the identity is over-identified, and the comparison of $\nFBI$ with $\widehat N_V\hat\rho_{\mathrm{pol}}\hat\rho_{\mathrm{cls}}\hat\rho_{\mathrm{sub}}$ enters the joint model as a likelihood term rather than a constraint.

\begin{table}[t]
\centering
\caption{What identifies each factor of the funnel (\cref{prop:funnel}(c)), and what does not.}\label{tab:funnel-sources}
\small
\begin{tabularx}{\textwidth}{@{}l>{\raggedright\arraybackslash}X l>{\raggedright\arraybackslash}X@{}}
\toprule
Factor & Source & Assumption & Not identified without further input \\
\midrule
$N_V$ & MSE on the linked lists (\cref{sec:identification,sec:models}) & \ref{asm:funnel-mse}, \cref{asm:fw-dependence} & the top interaction $\omega$ (\cref{thm:nonid}); a sensitivity range for it \\
$\rvic$ & victim survey: was the incident reported, by the respondent or anyone else, to the police or to an organisation that records incidents? & \ref{asm:funnel-survey} & channel self-detection; incidents outside the survey frame (institutional targets, non-household victims); state-level variation \\
$\rpol$ & victim survey: share reported to the police among incidents reported to some channel & \ref{asm:funnel-survey} & state-level variation; rare-event sampling error, so a prior is needed \\
$\rcls$ & police--NGO overlap among incidents that NGOs code as antisemitic and flag as reported to the police & \ref{asm:funnel-overlap}, \ref{asm:funnel-part}, linkage correction & without the police-report flag, only $\rpol\rcls\rsub$ among NGO-recorded incidents \\
$\rsub$ & agency participation data & \ref{asm:funnel-part} & incidence-weighted coverage; prior on $\sigma_c$ (\cref{sec:models-sae}) \\
\bottomrule
\end{tabularx}
\end{table}

\paragraph{What is not identified} Five gaps remain, to be closed by priors, external data or a reported range.
\begin{enumerate}[(1)]
  \item \emph{The unrecorded.} An incident that reaches no recording channel cannot appear on any list, however good the linkage, so every MSE estimator in this paper estimates at most $N_V$, and only under \ref{asm:funnel-mse}. Recovering $N=N_V/\rvic$ needs a victim survey. Without one we report $N_V$ and the downstream stages and treat $\rvic$ as a sensitivity parameter. Such incidents change only the unobservable cell: in expectation they add $N(1-\rvic)$ to $\mu_{\zero}$ and leave $\bm\mu_+$ unchanged, so by \eqref{eq:fw-key} targeting $N$ rather than $N_V$ is equivalent to raising $\bar\omega$ by
  \begin{equation}\label{eq:funnel-vic}
    \Delta_{\mathrm{vic}}=\log\Bigl\{1+\frac{(1-\rvic)\,N_V}{\rvic\,\mu_{\zero}}\Bigr\}\ge0,
  \end{equation}
  where $\mu_{\zero}$ is the expected number of unrecorded incidents in $U_d^V$. The $\bar\omega$ sensitivity analysis therefore covers $\rvic$ once its upper end is extended by $\Delta_{\mathrm{vic}}$ (\cref{sec:plan-sensitivity}). \ref{asm:funnel-mse} also demands that the dependence restriction hold on $U_d^V$. A list that detects incidents itself, such as a security patrol, enlarges $U_d^V$, but type-dependent detection can break the restriction even when every probability is positive: with two equally common types, a patrol list detecting with probabilities $0.6$ and $0.05$ and an NGO list with $0.5$ and $0.05$, independent given type, Lincoln--Petersen recovers $59\%$ of the truth in expectation. Detection also escapes surveys, which is why \ref{asm:funnel-survey} rules it out within the survey frame; where it occurs there, the survey understates $\rvic$ and $\widehat N_V/\hat\rho_{\mathrm{vic}}$ overstates $N$.
  \item \emph{Survey scope.} Household surveys of the NCVS type cover persons aged 12 or older in sampled households \citep{masucci2017,bjs2021hate}, so they say nothing about $\rvic$ or $\rpol$ for incidents whose only victim is an institution, such as a synagogue or cemetery; these need priors, informed where possible by surveys of institutions.
  \item \emph{Separating $\rpol$ from $\rcls$.} Incidents reported to the police but not classified as antisemitic are invisible in police data and appear only in NGO records that note a police report. Without that flag, the overlap identifies only the product of the downstream stages among NGO-recorded incidents, which also differs from the population product under referral (\cref{prop:referral}).
  \item \emph{Failures of \ref{asm:funnel-overlap}.} If an NGO learns of incidents from police statements, NGO capture depends on classification and the overlap overstates $\rcls$. Linkage errors distort the overlap in both directions (under LE a true pair survives with probability $(1-\delta)^2$), so the overlap is computed from the cut posterior of \cref{sec:models-link}.
  \item \emph{Participation is not submission of every incident.} \ref{asm:funnel-part} treats an agency that reports zero hate crimes as participating, as the FBI's counts do: for 2024, 19.0\% of participating agencies reported at least one hate-crime incident \citep{fbi2025summary}. An agency that never classifies anything as a bias crime contributes to $c_d$ but not to $\nFBI$, a loss correctly attributed to $\rcls$ only if $\rcls$ is measured.
\end{enumerate}
Finally, national surveys rarely support state-level estimates of $\rvic$ and $\rpol$, so in the cross-jurisdiction analysis below their variation comes from a prior or a residual. We then pool unmeasured stages into a single residual player and say so.

\subsection{Attributing differences to stages}\label{sec:funnel-shapley}

Consider jurisdictions $s=1,\dots,S$ with weights $w_s>0$, $\sum_sw_s=1$ (population shares or equal weights), and recorded rates $r_s=n_{\mathrm{FBI},s}/P_s$. By \cref{prop:funnel}(a), in its realised or its expected form, $r_s=\theta_s\,\rho_{\mathrm{vic},s}\rho_{\mathrm{pol},s}\rho_{\mathrm{cls},s}\rho_{\mathrm{sub},s}$ with $\theta_s=N_s/P_s$. Taking logs,
\begin{equation}\label{eq:funnel-logs}
  Y_s=\log r_s=\sum_{j=1}^{m}X_{js},\qquad
  (X_{1s},\dots,X_{5s})=(\log\theta_s,\log\rho_{\mathrm{vic},s},\log\rho_{\mathrm{pol},s},\log\rho_{\mathrm{cls},s},\log\rho_{\mathrm{sub},s}),
\end{equation}
with $m=5$ players: incidence and the four stages. Let $\Var_w$ and $\Cov_w$ denote weighted moments across jurisdictions. The Shapley value \citep{shapley1953} of a game $v$ on the players $[m]=\{1,\dots,m\}$, with $v(\emptyset)=0$, allocates the worth $v([m])$ as
\begin{equation}\label{eq:funnel-shapley-def}
  \varphi_j(v)=\sum_{A\subseteq[m]\setminus\{j\}}\frac{|A|!\,(m-1-|A|)!}{m!}\,\{v(A\cup\{j\})-v(A)\},
\end{equation}
the average marginal contribution of $j$ over all orderings of the players. Following \citet{shorrocks2013}, a factor is ``eliminated'' by holding it at a reference value. For dispersion the reference is the cross-jurisdiction mean; for a comparison of two jurisdictions it is the first jurisdiction's value.

\begin{theorem}[Shapley decomposition of the funnel]\label{thm:shapley}
\leavevmode
\begin{enumerate}[(i)]
  \item \emph{Dispersion.} For the game $v(A)=\Var_w\big(\sum_{j\in A}X_j\big)$, the Shapley value is
  \[
    \varphi_j=\Cov_w(X_j,Y)=\Var_w(X_j)+\sum_{i\ne j}\Cov_w(X_i,X_j).
  \]
  The shares sum to $\Var_w(Y)$, and $\varphi_j<0$ if and only if $\Cov_w(X_j,Y)<0$.
  \item \emph{Two jurisdictions, levels.} Let jurisdictions $A$ and $B$ have factor vectors $\bm a,\bm b\in[0,\infty)^m$, and $r_A=\prod_ja_j$, $r_B=\prod_jb_j$. For the switching game $v(T)=\prod_{j\in T}b_j\prod_{j\notin T}a_j-r_A$, the Shapley value is
  \begin{equation}\label{eq:funnel-shapley-levels}
  \begin{aligned}
    \varphi_j&=(b_j-a_j)\sum_{T\subseteq[m]\setminus\{j\}}\frac{|T|!\,(m-1-|T|)!}{m!}\prod_{i\in T}b_i\prod_{i\notin T\cup\{j\}}a_i\\
    &=(b_j-a_j)\int_0^1\prod_{i\ne j}\{a_i+t(b_i-a_i)\}\,dt .
  \end{aligned}
  \end{equation}
  It is exact, $\sum_j\varphi_j=r_B-r_A$. It is antisymmetric under exchanging $A$ and $B$, and $\varphi_j=0$ when $a_j=b_j$. For strictly positive factors, the LMDI-I allocation $\varphi_j^{\mathrm L}=\mathcal L(r_A,r_B)\log(b_j/a_j)$, with the logarithmic mean $\mathcal L(x,y)=(x-y)/(\log x-\log y)$ and $\mathcal L(x,x)=x$, is also exact. Both allocations are line integrals of $\nabla f$, where $f(\bm x)=\prod_ix_i$, from $\bm a$ to $\bm b$. Shapley integrates along the segment $\bm a+t(\bm b-\bm a)$, and LMDI along the geometric path $\bm a^{1-t}\bm b^{t}$. With $d_i=\log(b_i/a_i)$,
  \begin{equation}\label{eq:funnel-cubic}
    \varphi_j-\varphi_j^{\mathrm L}=\frac{\sqrt{r_Ar_B}}{12}\,d_j\sum_{i=1}^{m}d_i(d_i-d_j)+O(\|\bm d\|^5).
  \end{equation}
  \item \emph{Posterior uncertainty.} Let the factors be uncertain, with posterior means $\mu_{js}=\E[X_{js}\mid\mathrm{data}]$ and $\mu_{Ys}=\E[Y_s\mid\mathrm{data}]$, and write $\bar X_j=\sum_sw_sX_{js}$ and $\bar Y=\sum_sw_sY_s$. The posterior of $\varphi_j$ in (i) is obtained by evaluating $\Cov_w(X_j,Y)$ at each posterior draw, and
  \begin{equation}\label{eq:funnel-postmean}
    \E[\varphi_j\mid\mathrm{data}]=\Cov_w(\mu_j,\mu_Y)+\sum_sw_s\Cov(X_{js},Y_s\mid\mathrm{data})-\Cov(\bar X_j,\bar Y\mid\mathrm{data}).
  \end{equation}
  Summed over $j$, the correction terms equal $\sum_sw_s\Var(Y_s\mid\mathrm{data})-\Var(\bar Y\mid\mathrm{data})\ge0$. If $Y$ is observed (realised recorded rates), the correction vanishes, and the plug-in share $\Cov_w(\mu_j,Y)$ is the posterior mean.
\end{enumerate}
\end{theorem}

\begin{proof}
(i) Let $\Sigma$ be the weighted covariance matrix of $(X_1,\dots,X_m)$. Then $v(A)=\sum_{i,k\in A}\Sigma_{ik}$, and for $j\notin A$ the marginal contribution is $v(A\cup\{j\})-v(A)=\Sigma_{jj}+2\sum_{i\in A}\Sigma_{ij}$. In a uniformly random ordering, each $i\neq j$ precedes $j$ with probability $1/2$. Averaging over orderings therefore gives $\varphi_j=\Sigma_{jj}+\sum_{i\ne j}\Sigma_{ij}=\Cov_w(X_j,\sum_iX_i)$. Summing over $j$ gives $\sum_{i,j}\Sigma_{ij}=\Var_w(Y)$.

(ii) Since $v(T\cup\{j\})-v(T)=(b_j-a_j)\prod_{i\in T}b_i\prod_{i\notin T\cup\{j\}}a_i$, the first expression in \eqref{eq:funnel-shapley-levels} is \eqref{eq:funnel-shapley-def}. For the second, expand $\prod_{i\ne j}\{(1-t)a_i+tb_i\}=\sum_{T}t^{|T|}(1-t)^{m-1-|T|}\prod_{i\in T}b_i\prod_{i\notin T\cup\{j\}}a_i$ and use $\int_0^1t^{q}(1-t)^{m-1-q}dt=q!(m-1-q)!/m!$. Exactness is the efficiency of the Shapley value: $\sum_j\varphi_j=v([m])=r_B-r_A$. Antisymmetry follows from the substitution $t\mapsto1-t$. Along $\bm x(t)=\bm a+t(\bm b-\bm a)$, $\partial_jf(\bm x(t))\,\dot x_j(t)=(b_j-a_j)\prod_{i\ne j}x_i(t)$, which is the integrand in \eqref{eq:funnel-shapley-levels}. Along $\bm x(t)=\bm a^{1-t}\bm b^{t}$, $\partial_jf(\bm x(t))\,\dot x_j(t)=f(\bm x(t))\,d_j=r_A^{1-t}r_B^{t}d_j$, whose integral over $[0,1]$ is $\mathcal L(r_A,r_B)\,d_j=\varphi_j^{\mathrm L}$. Exactness of both allocations is the gradient theorem, $\sum_j\int\partial_jf\,dx_j=f(\bm b)-f(\bm a)$; for LMDI it can also be checked directly, since $\mathcal L(r_A,r_B)\sum_jd_j=\mathcal L(r_A,r_B)\log(r_B/r_A)=r_B-r_A$. The expansion \eqref{eq:funnel-cubic} is proved in \cref{app:bayes-proofs}.

(iii) Each $\varphi_j$ is a function of the factors, so its posterior is the image of their posterior. Write $\varphi_j=\sum_sw_sX_{js}Y_s-\bar X_j\bar Y$ and take posterior expectations, using $\E[UV]=\E U\,\E V+\Cov(U,V)$ for each product. Summing over $j$ gives $\sum_sw_s\Var(Y_s\mid\mathrm{data})-\Var(\bar Y\mid\mathrm{data})$. This is nonnegative, because $\Var(\sum_sw_sY_s)\le(\sum_sw_s\sigma_s)^2\le\sum_sw_s\sigma_s^2$ with $\sigma_s^2=\Var(Y_s\mid\mathrm{data})$, by the Cauchy--Schwarz and Jensen inequalities. If $Y$ is observed, $\varphi_j$ is linear in the unknowns, and every correction term is zero.
\end{proof}

Part (i) is elementary: the covariance shares $\Cov_w(X_j,Y)/\Var_w(Y)$ are the ``natural'' decomposition of a variance by factor components \citep{shorrocks1982}, and the proof shows that they are also the Shapley value of the elimination game. The closed form in (ii) is an instance of Owen's multilinear extension \citep{owen1972}. Shapley decompositions of multiplicative identities are standard in index decomposition analysis \citep{albrecht2002}, as is LMDI \citep{ang2004,ang2015}; to our knowledge, the cubic expansion \eqref{eq:funnel-cubic} is not stated elsewhere.
\emph{Both} decompositions are exact, with no residual; they are different allocations of the same total, agreeing when all factors differ by the same proportion and to second order in general, and differing materially when a factor differs twofold or more between jurisdictions, which is common for participation. LMDI's shares equal the log shares $d_j/\sum_id_i$, consistent with part (i), but LMDI is undefined when a factor is zero (for example, no participating agency), whereas \eqref{eq:funnel-shapley-levels} is not. We report the Shapley allocation for levels and LMDI as a check; with more than two jurisdictions, (ii) is applied to each against a common reference such as the national factors.

\paragraph{Reading the shares} A positive $\varphi_j$ in (i) means that stage $j$ widens the spread of recorded rates; a negative share means that factor $j$ is negatively correlated with the recorded rate across jurisdictions. The most consequential case is incidence. If $\varphi_1<0$, recorded rates tend to \emph{fall} as incidence rises. In that case, places with more antisemitic incidents record fewer per capita, because their funnel loses more, and a map of recorded rates tends to rank jurisdictions in the opposite order to incidence. Restricting the game to the four stages decomposes instead the dispersion of $\log(r_s/\theta_s)=-\log D_{1,s}$. By part (iii), shares are computed draw by draw: plugging in posterior means understates dispersion when $Y$ is latent, because posterior means are shrunk toward the national mean \citep{louis1984}. We decompose realised rates where counts are large, and expected rates where counts are small or zero.

\subsection{An illustration with invented numbers}\label{sec:funnel-illustration}

\Cref{tab:funnel-illus} applies \cref{thm:shapley}(ii) to two hypothetical jurisdictions. \textbf{All inputs are invented for exposition} and describe no real place. Jurisdiction $B$ has $50\%$ more incidents per capita than $A$, yet its recorded rate is less than a third of $A$'s. Incidence contributes $+0.390$ per $100{,}000$ to the difference $r_B-r_A=-0.974$, that is $-40\%$ of the gap. Classification ($51\%$) and submission ($58\%$) together account for more than the whole gap, police reporting for $23\%$ and reaching a channel for $8\%$; with $A$'s classification and submission rates, $B$'s recorded rate would be $1.536$, above $A$'s. The LMDI shares differ from the Shapley shares by $0.5$ to $4.3$ percentage points; the largest difference, for incidence, is $0.042$ per $100{,}000$, and the cubic term of \eqref{eq:funnel-cubic} gives $0.039$.

\begin{table}[t]
\centering
\caption{Shapley and LMDI decompositions of the difference in recorded rates between two \emph{hypothetical} jurisdictions (\cref{thm:shapley}(ii)). \textbf{Illustrative, invented numbers; not data.} Rates are per $100{,}000$ residents; percentages are shares of $r_B-r_A$ and may not sum to $100$ because of rounding.}\label{tab:funnel-illus}
\small
\begin{tabular}{@{}lccc rr rr@{}}
\toprule
 & & & & \multicolumn{2}{c}{Shapley} & \multicolumn{2}{c}{LMDI-I}\\
\cmidrule(lr){5-6}\cmidrule(lr){7-8}
Factor & $A$ & $B$ & $\log(b_j/a_j)$ & $\varphi_j$ & \% & $\varphi^{\mathrm L}_j$ & \% \\
\midrule
Incidence $\theta$       & 8.00 & 12.00 & $+0.405$ & $+0.390$ & $-40.0$ & $+0.348$ & $-35.7$ \\
Reaching a channel $\rvic$ & 0.60 & 0.55  & $-0.087$ & $-0.080$ & $8.2$   & $-0.075$ & $7.7$ \\
Police report $\rpol$    & 0.45 & 0.35  & $-0.251$ & $-0.228$ & $23.4$  & $-0.216$ & $22.2$ \\
Classification $\rcls$   & 0.70 & 0.40  & $-0.560$ & $-0.494$ & $50.7$  & $-0.481$ & $49.3$ \\
Submission $\rsub$       & 0.95 & 0.50  & $-0.642$ & $-0.562$ & $57.7$  & $-0.551$ & $56.6$ \\
\midrule
Recorded rate $r$        & 1.436 & 0.462 & $-1.134$ & $-0.974$ & $100$ & $-0.974$ & $100$ \\
\bottomrule
\end{tabular}
\end{table}

\subsection{Policy reading}\label{sec:funnel-policy}

Each stage corresponds to a different lever.
\begin{description}[leftmargin=1.2em,style=nextline]
  \item[Incidence ($\theta$).] Prevention, security and education; not a reporting problem. A high recorded rate may reflect a well-functioning funnel, and a low one may hide high incidence.
  \item[Reaching a channel ($\rvic$).] Victim support and community trust, which raise disclosure, and accessible channels such as multilingual, anonymous and third-party reporting.
  \item[Police reporting ($\rpol$).] Accessible and trusted police: liaison officers, online reporting, and forwarding of reports by community organisations with the victim's consent.
  \item[Classification ($\rcls$).] Classification protocols and training, documented bias indicators, and periodic audits comparing police and NGO coding of the same incidents (\cref{prop:funnel}(c)).
  \item[Submission ($\rsub$).] NIBRS participation and data quality: technical assistance, funding for records systems, and follow-up on agencies that persistently report zero.
\end{description}
Three cautions apply. Shapley shares account for observed differences; they are not causal effects of interventions, and a change that raises trust in the police may raise $\rvic$ as well as $\rpol$. Improving the funnel raises recorded counts, and applied to two years instead of two places, \cref{thm:shapley}(ii) separates a rise due to better reporting from a rise in incidence, a separation needed before a recorded increase is read as an increase in antisemitism. And a share is only as informative as the data behind that stage's variation; shares of prior-driven stages are reported as such or pooled into a residual.
